%% 
%% Copyright 2019-2024 Elsevier Ltd
%% 
%% This file is part of the 'CAS Bundle'.
%% --------------------------------------
%% 
%% It may be distributed under the conditions of the LaTeX Project Public
%% License, either version 1.3c of this license or (at your option) any
%% later version.  The latest version of this license is in
%%    http://www.latex-project.org/lppl.txt
%% and version 1.3c or later is part of all distributions of LaTeX
%% version 1999/12/01 or later.
%% 
%% The list of all files belonging to the 'CAS Bundle' is
%% given in the file `manifest.txt'.
%% 
%% Template article for cas-sc documentclass for 
%% double column output.

\documentclass[a4paper,fleqn]{cas-sc}

% If the frontmatter runs over more than one page
% use the longmktitle option.

%\documentclass[a4paper,fleqn,longmktitle]{cas-sc}

%\usepackage[numbers]{natbib}
%\usepackage[authoryear]{natbib}
\usepackage[authoryear,longnamesfirst]{natbib}

%%%Author macros
\def\tsc#1{\csdef{#1}{\textsc{\lowercase{#1}}\xspace}}
\tsc{WGM}
\tsc{QE}
%%%

% Uncomment and use as if needed
%\newtheorem{theorem}{Theorem}
%\newtheorem{lemma}[theorem]{Lemma}
%\newdefinition{rmk}{Remark}
%\newproof{pf}{Proof}
%\newproof{pot}{Proof of Theorem \ref{thm}}

\begin{document}
\let\WriteBookmarks\relax
\def\floatpagepagefraction{1}
\def\textpagefraction{.001}

% Short title
\shorttitle{ML-Based Football Squad Optimization}    

% Short author
\shortauthors{Nunes and Pereira}  

% Main title of the paper
\title [mode = title]{Machine Learning-Based Optimization of Football Squad Composition Under Budget Constraints}  

% Title footnote mark
% eg: \tnotemark[1]
\tnotemark[1] 

% Title footnote 1.
% eg: \tnotetext[1]{Title footnote text}
\tnotetext[1]{This study was financed in part by the Coordenação de Aperfeiçoamento de Pessoal de Nível Superior – Brasil (CAPES) – Finance Code 001.}

% First author
%
% Options: Use if required
% eg: \author[1,3]{Author Name}[type=editor,
%       style=chinese,
%       auid=000,
%       bioid=1,
%       prefix=Sir,
%       orcid=0000-0000-0000-0000,
%       facebook=<facebook id>,
%       twitter=<twitter id>,
%       linkedin=<linkedin id>,
%       gplus=<gplus id>]

\author[1]{Lucio Vargas de Albuquerque Nunes\corref{cor1}}

\ead{luciovanunes@gmail.com}
\ead[url]{https://orcid.org/0000-0001-8803-3293}

\author[1]{Dilson Lucas Pereira}

\ead{dilson.pereira@ufla.br}
\ead[url]{https://orcid.org/0000-0002-6307-5152}

\cortext[cor1]{Corresponding author}

\affiliation[1]{
organization={Universidade Federal de Lavras},
city={Lavras},
state={MG},
country={Brazil}
}

% Here goes the abstract
\begin{abstract}
Football is characterized by high levels of uncertainty, which makes squad planning and player recruitment complex decision-making processes. 
This study investigates the use of Machine Learning (ML) and Operations Research (OR) techniques as decision-support tools for football squad composition. 
A predictive model based on a multilayer perceptron (MLP) is trained to estimate the final number of points obtained by football teams in a season using 
individual player statistics from the previous season. 
The predicted team performance is then used as the objective function in heuristic optimization algorithms that simulate player recruitment under budget constraints. 
Experiments conducted using data from major European leagues demonstrate the potential of combining predictive modeling and 
optimization to support squad planning decisions. The proposed framework provides a structured and data-driven approach to evaluating recruitment strategies, 
reducing subjectivity in player acquisition and contributing to more informed decision-making in professional football.
\end{abstract}

% Use if graphical abstract is present
%\begin{graphicalabstract}
%\includegraphics{}
%\end{graphicalabstract}

% Research highlights
\begin{highlights}
\item Machine learning predicts team performance from player statistics ($R^2 \approx 0.55$)
\item Greedy and local search algorithms optimize squad composition under budget constraints
\item Framework simulates realistic football transfer scenarios
\item Experiments conducted using Premier League and Serie A datasets
\end{highlights}

% Keywords
% Each keyword is seperated by \sep
\begin{keywords}
sports analytics \sep football \sep machine learning \sep squad optimization \sep decision support
\end{keywords}

\maketitle

% Main text
\section{Introduction}

Football is the most popular sport in the world \cite{Almulla2023}. Over the past decades, the sport has undergone significant technological and analytical 
development, accompanied by substantial economic growth in the global football industry \cite{Rodrigues2022}. Despite these advances, football remains highly 
unpredictable. A notable example is Leicester City's unexpected triumph in the Premier League (PL) in 2016, an outcome that surprised analysts 
and fans worldwide \cite{Baboota2019}.

Given the economic relevance and global popularity of football, both academic researchers and industry professionals have increasingly explored 
data-driven approaches to support decision-making in the sport \cite{Baboota2019}. The financial magnitude of football is reflected in the transfer market, 
which moves billions of dollars annually. For instance, between June and September 2023, global player transfers reached approximately USD 7.3 billion according 
to FIFA reports \cite{Globo2023}. These large financial investments highlight the importance of improving decision-making processes related to player recruitment 
and squad construction.

In the field of Computer Science, several studies have investigated the use of statistical analysis and machine learning (ML) techniques to model 
patterns and generate predictions in complex environments. Joseph, Fenton and Neil (2006), for example, investigated the use of Bayesian networks and 
other ML techniques to predict match outcomes for Tottenham Hotspur during the 1995–1997 seasons. In professional football practice, data-driven strategies 
have also been adopted by clubs such as Brighton \& Hove Albion, which has implemented a recruitment model strongly supported by statistical analysis. This 
approach has enabled the identification of undervalued players in the transfer market, generating substantial financial returns through successful 
transfers \cite{BBC2023Caicedo, CNN2023Caicedo, ESPN2023}.

In academic research, most studies applying ML to football focus primarily on predicting match outcomes, estimating win probabilities, or optimizing 
betting strategies. However, comparatively fewer studies address the problem of squad composition and player recruitment using data-driven 
approaches \cite{Abidin2021}. To the best of our knowledge, few studies attempt to predict season-level team performance from individual player statistics 
and incorporate these predictions into a formal squad optimization framework. This gap is particularly relevant because squad construction requires 
evaluating multiple technical, physical and contextual factors simultaneously, including team playing style, league characteristics and player career 
stage \cite{Abidin2021}. Furthermore, the stochastic nature of football performance means that even seemingly optimal transfers may fail to 
deliver the expected results.

Transfer market inefficiencies further illustrate the complexity of these decisions. According to Graham (2024), a transfer can be considered 
successful if the player starts at least 50\% of matches during the first two seasons with the club, a concept referred to as the ``50\% rule''. Based on 
this criterion, approximately 46\% of transfers above \EUR{10 million} in the EPL between 1992 and 2021 can be classified as unsuccessful \cite{Graham2024}. These 
outcomes are often associated with factors such as tactical incompatibility, inaccurate evaluation of player ability, positional mismatches, coaching 
decisions or physical issues.

Recent advances in ML have significantly expanded the capacity to analyze sports data, enabling the identification of patterns that were previously 
difficult to detect in large datasets. However, predictive models do not eliminate the inherent uncertainty of football; rather, they serve as decision-support 
tools that can help reduce risk and support strategic choices.

To address this challenge, this study proposes an integrated framework that combines machine learning-based performance prediction with heuristic squad optimization.

The main contributions of this work are threefold. First, we develop a machine learning model capable of predicting season-level team performance 
from individual player statistics. Second, we integrate these predictions into an optimization framework that simulates squad recruitment decisions under 
budget constraints. Third, we demonstrate through computational experiments that the proposed approach can identify recruitment strategies capable of improving 
predicted team performance.

\section{Related Work}

\section{Data and Feature Engineering}

\section{Machine Learning Model}

\section{Squad Optimization Framework}

\section{Experimental Results}

\section{Conclusion}

% Numbered list
% Use the style of numbering in square brackets.
% If nothing is used, default style will be taken.
%\begin{enumerate}[a)]
%\item 
%\item 
%\item 
%\end{enumerate}  

% Unnumbered list
%\begin{itemize}
%\item 
%\item 
%\item 
%\end{itemize}  

% Description list
%\begin{description}
%\item[]
%\item[] 
%\item[] 
%\end{description}  

\clearpage %%Remove this from your manuscript


% Figure
\begin{figure}%[]
  \centering
%    \includegraphics{}
    \caption{}\label{fig1}
\end{figure}


\begin{table}%[]
\caption{}\label{tbl1}
\begin{tabular*}{\tblwidth}{@{}LL@{}}
\toprule
  &  \\ % Table header row
\midrule
 & \\
 & \\
 & \\
 & \\
\bottomrule
\end{tabular*}
\end{table}

% Uncomment and use as the case may be
%\begin{theorem} 
%\end{theorem}

% Uncomment and use as the case may be
%\begin{lemma} 
%\end{lemma}

%% The Appendices part is started with the command \appendix;
%% appendix sections are then done as normal sections
%% \appendix

\section{}\label{}

% To print the credit authorship contribution details
\printcredits

%% Loading bibliography style file
%\bibliographystyle{model1-num-names}
\bibliographystyle{cas-model2-names}

% Loading bibliography database
\bibliography{cas-refs}

% Biography
%\bio{}
% Here goes the biography details.
%\endbio

%\bio{pic1}
% Here goes the biography details.
%\endbio

\end{document}

